\section{Geometric certificates for the core and its marked attachments}
\label{geo:section}

A bounded complementary core is useful only if its topology can be
certified and the omitted product regions can be glued back correctly.
This section supplies a conservative consumer of the existing native
normal-cut geometry.  The consumer has two parts.  First, it certifies
handlebodies by cancellations witnessed by \emph{actual geometric
intersection counts}.  Second, it reconstructs the topology of an assembly
whose interfaces are separated disks.  A complete polynomial criterion is
obtained for the ball cores accessible to raw singleton cancellation.
An explicit infinite family shows that this criterion can succeed uniformly
when the normal coordinates grow exponentially.

The algebraic source-closure principle used below is already present in
ProveIt report~56, \texttt{reports/56/article/raw\_epoch.tex}
\cite{ProveIt2026}.  We prove the needed matrix statement for completeness.
The contribution here is its geometric realization and its application to
source-authenticated normal-cut handles.  Neither the classical handle
cancellation theorem nor the elementary topology of interval bundles is
claimed as new.

\subsection{The geometric input and the role of literal occurrences}
\label{geo:input}

We use the full geometric conclusion of the existing source-parallelity
construction, described in \texttt{synthesis/parallelity.tex}
\cite{ProveIt2026}.  This conclusion supplies an embedded interval bundle,
its complete base polygons and fiber identifications, and an inherited
complementary handle structure.  The polygon construction includes contacts
of dimensions zero--one, zero--two, and one--two, together with the source
edge and vertex strata.  Its core records contain both ends of every
one-handle, the alternating attaching sectors of every two-handle, and the
vertex-link charts of every three-handle.  We invoke this geometric
conclusion, rather than infer an embedding from Euler characteristics,
component histograms, or an abstract presentation.

\begin{definition}[Authenticated core incidence data]
\label{geo:def-incidence}
Let $K$ be a compact orientable three-manifold with a genuine finite handle
decomposition.  Suppose that it has no three-handles.  Its zero- and
one-handles form a disjoint union of handlebodies.  Contract a spanning
forest of their incidence graph by geometric zero/one-handle cancellations.
In one resulting connected component, let
\[
 D_1,\ldots,D_k
\]
be the complete system of co-core disks of the remaining one-handles, and
let $\alpha_1,\ldots,\alpha_f$ be the disjoint embedded two-handle attaching
curves on the boundary of this handlebody.  Put the curves and disk
boundaries in the transverse position specified by the handle charts and
set
\[
 n_{ih}=\bigl|\alpha_i\cap\partial D_h\bigr|.
 \tag{G1}\label{geo:eq-incidence}
\]
The matrix $(n_{ih})$ is authenticated when every counted occurrence is an
actual passage through the corresponding source one-handle, with its
source attaching sector retained.
\end{definition}

A forest cancellation only absorbs tree one-handles into zero-handles.
The other co-core disks survive, so deleting the tree-handle entries from
the source attaching lists gives the incidences in
\eqref{geo:eq-incidence}.  A two-handle cannot connect distinct components
of the zero/one-handle union: its attaching circle is connected.  Thus the
forest also identifies the components relevant to all subsequent
cancellations.

This contract is stronger than a cellular boundary matrix or a
fundamental-group presentation.  An exponent sum of $1$, or a letter
occurring once after free reduction, need not mean that an attaching curve
meets a belt circle once.  All occurrences in
\eqref{geo:eq-incidence} are unsigned \emph{geometric} occurrences.  There
is no free reduction, cancellation of opposite signs, or replacement by a
homotopic curve hidden in this definition.

\subsection{A geometric realization of positive incidence elimination}
\label{geo:slides}

\begin{lemma}[Meridian-disk realization of a raw singleton pivot]
\label{geo:lem-slide}
Let $H$ be an orientable handlebody with a complete disk system
$D_1,\ldots,D_k$, and let $\alpha_1,\ldots,\alpha_f$ be pairwise disjoint
embedded attaching curves on $\partial H$.  Write
$n_{ih}=|\alpha_i\cap\partial D_h|$.  Suppose that $n_{rg}=1$.
Attach the two-handle with attaching curve $\alpha_r$.  The resulting
manifold is a handlebody of genus $k-1$ with a complete disk system indexed
by $h\ne g$ for which the surviving attaching curves have incidences
\[
 n'_{ih}=n_{ih}+n_{ig}n_{rh},
 \qquad i\ne r,\quad h\ne g.
 \tag{G2}\label{geo:eq-positive-pivot}
\]
In particular, the equality counts actual intersections; it is not merely
an equality of homology classes or group words.
\end{lemma}

\begin{proof}
Keep $D_g$ fixed, and write $q$ for its unique intersection with
$\alpha_r$.  If $\alpha_r$ meets another disk boundary, choose such an
intersection $p$ adjacent to $q$ along $\alpha_r$, among all current
intersections with disk boundaries.  Suppose that $p$ lies on
$\partial D_h$.  The intervening arc has interior disjoint from all disk
boundaries.  It is also disjoint from every other attaching curve, since
those curves are disjoint from $\alpha_r$.

Slide $D_h$ over a parallel copy of $D_g$ along a thin band following this
arc.  Choose the copy on the side of $D_g$ approached by the arc, and choose
its endpoint just before $q$.  The band then misses $D_g$ itself.  The
standard boundary connected sum replaces $D_h$ by a properly embedded disk
and preserves a disjoint complete disk system.  This elementary disk slide
can be seen in a regular neighborhood of the two disks and the band:
sliding one disk across the parallel copy changes the chosen meridian
system without changing the handlebody cut out by it.

The smoothing at the band removes the intersection at $p$ and the unique
intersection of $\alpha_r$ with the parallel copy of $\partial D_g$.
The two sides of the band miss $\alpha_r$.  Thus the replacement for $D_h$
has exactly one fewer intersection with $\alpha_r$; the intersections of
$\alpha_r$ with the other disks are unchanged.  Small band endpoints and a
sufficiently thin band miss all other attaching curves.  For any $i\ne r$,
the new disk boundary consequently consists, for purposes of counting
intersections with $\alpha_i$, of the old $\partial D_h$ and one parallel
copy of $\partial D_g$.  Its intersection count increases by exactly
$n_{ig}$.

Repeat this operation, always using the fixed disk $D_g$ and an adjacent
intersection along $\alpha_r$.  At each step the total number of
intersections of $\alpha_r$ with the other disk boundaries decreases by
one.  Hence the process terminates.  The disk indexed by $h$ has been slid
exactly $n_{rh}$ times, whereas $D_g$ and its intersections with the other
attaching curves have remained fixed.  The resulting disks $E_h$, $h\ne
g$, are disjoint from $\alpha_r$ and satisfy
\[
 |\alpha_i\cap\partial E_h|=n_{ih}+n_{ig}n_{rh}.
\]

In the disk system consisting of $D_g$ and the $E_h$, the selected
attaching curve goes over the handle dual to $D_g$ exactly once and misses
all other co-cores.  Its one/two-handle pair therefore cancels by the
geometric one-point cancellation theorem; see, for example,
\cite[Section~10.1.3, Lemma~10.3]{SadykovNotes}.  Equivalently, cutting along
the $E_h$ leaves a solid torus, and attaching a two-handle along a curve
meeting its meridian once turns that solid torus into a ball.  Thus the
$E_h$ form a complete disk system after attachment.  Every other attaching
curve is disjoint from $\alpha_r$, so it lies on the new boundary, and its
intersections with the $E_h$ are unchanged.  This proves
\eqref{geo:eq-positive-pivot}.
\end{proof}

\begin{remark}[What the construction does and does not record]
\label{geo:rem-geometric-existence}
The proof permits very many disk slides.  Its purpose is to prove that a
specific integer update is geometrically realizable.  A classifier can
verify that update without expanding the disks or constructing the
resulting boundary homeomorphism.  If a later hierarchy operation needs an
explicit lifted disk or the transported position of a peripheral curve,
that additional output must have its own representation and complexity
bound.  The present certificate asserts the handlebody type and the
existence of the indicated complete disk system.
\end{remark}

\begin{proposition}[Saturation removes occurrence growth from verification]
\label{geo:prop-saturation}
Let $C=\{0,1,2\}$ have operations
\[
 a\mathbin{\oplus}b=\min\{2,a+b\},\qquad
 a\mathbin{\otimes}b=\min\{2,ab\},
\]
and let $\sigma(n)=\min\{2,n\}$ for $n\in\mathbb N$.
Starting from authenticated incidences, store only $c_{ih}=\sigma(n_{ih})$.
A pivot $(r,g)$ is geometrically valid whenever $c_{rg}=1$.  After deleting
its row and column, the exact saturated incidences are
\[
 c'_{ih}
 =c_{ih}\mathbin{\oplus}(c_{ig}\mathbin{\otimes}c_{rh})
 =\min\{2,c_{ih}+c_{ig}c_{rh}\}.
 \tag{G3}\label{geo:eq-capped-pivot}
\]
For $k+f=O(T)$, a supplied sequence removing all two-handles is replayed in
$O(T^3)$ operations on constant-sized entries and $O(T^2)$ working space.
Its forest and pivot labels use $O(T\log(T+1))$ bits.
\end{proposition}

\begin{proof}
The map $\sigma$ preserves zero, one, addition and multiplication when its
codomain uses $\oplus$ and $\otimes$.  This follows directly by separating
the cases $0$, $1$, and at least $2$.  In particular,
$\sigma(n)=1$ if and only if $n=1$.  Lemma~\ref{geo:lem-slide} and induction
therefore show that every stored matrix is exactly the saturation of an
actual geometric incidence matrix.  No large integer or word needs to be
reconstructed.  There are at most $\min\{k,f\}$ pivots, each of which
updates at most $kf$ entries.  The remaining bounds follow from the matrix
size and from recording $O(T)$ source identifiers.
\end{proof}

The operation bound concerns this matrix phase after the source charts
have been authenticated.  It does not absorb the cost of input validation,
normal-coordinate arithmetic, the interval-orbit proofs used to obtain the
parallelity data, or a preceding search for the normal surface.

\begin{corollary}[Core handlebody certificate]
\label{geo:cor-core-handlebody}
Consider a connected authenticated core with $v$ zero-handles, $e$
one-handles, $f$ two-handles, and no three-handles.  If a sequence of valid
pivots removes every two-handle after spanning-tree contraction, the core
is a handlebody of genus
\[
 g=e-v+1-f.
 \tag{G4}\label{geo:eq-core-genus}
\]
\end{corollary}

\begin{proof}
Tree contraction leaves one zero-handle and $k=e-v+1$ one-handles.
Each pivot attaches one of the remaining two-handles and reduces the
handlebody genus by one, by Lemma~\ref{geo:lem-slide}.  After all $f$
attachments have been made, no further handles remain.  The result is a
handlebody of genus $k-f$.
\end{proof}

The hypothesis excluding three-handles is essential to this conclusion.
The actual attaching curves and their disjointness are essential as well:
an arbitrary integer matrix is not a certificate of a three-manifold's
topology.  Failure of a search for these pivots does not show that the core
is not a handlebody.

\subsection{A complete polynomial criterion within the ball-core class}
\label{geo:ball-closure}

There is a substantial improvement over trying arbitrary singleton
pivots.  In the square case, an entire successful raw elimination can be
discovered without forming any fill entries.  This is the geometric
application of the immutable-source theorem from report~56
\cite{ProveIt2026}.

\begin{lemma}[Reversing a successful positive unit pivot]
\label{geo:lem-permanent-pivot}
Let $A=(a_{ij})$ be a square matrix of nonnegative integers with $a_{dc}=1$.
For $i\ne d$ and $j\ne c$, put
\[
 b_{ij}=a_{ij}+a_{ic}a_{dj}.
\]
If $\operatorname{per}(B)=1$, then $\operatorname{per}(A)=1$.
The permanent of the empty matrix is understood to be $1$.
\end{lemma}

\begin{proof}
View a matrix entry as a number of separately labelled edges in a
bipartite multigraph.  A permanent counts its edge-labelled perfect
matchings.  Tag each edge of $B$ either as a direct edge of $A$, or as an
ordered pair consisting of an edge from $i$ to $c$ and an edge from $d$ to
$j$.

The unique tagged perfect matching of $B$ contains at most one pair-created
edge.  If it contained two, say in positions $(i,j)$ and $(k,\ell)$, then
$i\ne k$ and $j\ne\ell$.  Exchanging their destination factors would
produce pair-created edges in positions $(i,\ell)$ and $(k,j)$, giving a
second tagged perfect matching.  If the unique matching contains no
pair-created edge, adjoin the unique edge $(d,c)$ to obtain a perfect
matching of $A$.  If it contains one, replace that edge by its two
component edges to obtain a perfect matching of $A$.

Conversely, map any perfect matching of $A$ into a tagged perfect matching
of $B$.  If it uses $(d,c)$, delete that edge.  Otherwise there is a unique
row matched to $c$ and a unique column matched from $d$; replace these two
edges by the corresponding pair-created edge.  This map is injective:
the zero-or-one number of pair-created edges distinguishes the cases, and
the tags recover all original edge labels.  Hence $A$ has at least one and
at most one perfect matching.
\end{proof}

\begin{definition}[Source row-leaf deletion]
\label{geo:def-leaf}
In a matrix with some rows and columns already deleted, a row is a
\emph{source row leaf} if its immutable source entries have exactly one
nonzero among the live columns, and that entry equals $1$.  A row-leaf
deletion removes this row and its unique live column.  No fill update is
performed.
\end{definition}

\begin{theorem}[Complete source criterion for raw ball-core cancellation]
\label{geo:thm-ball-closure}
Let a connected authenticated core have no three-handles and satisfy
\[
 k=e-v+1=f.
\]
Let $A$ be its square unsigned incidence matrix after spanning-tree
contraction.  The following conditions are equivalent:
\begin{enumerate}
\item Some sequence of positive singleton pivots
      \eqref{geo:eq-positive-pivot} removes every row and column.
\item The source matrix satisfies $\operatorname{per}(A)=1$.
\item Source row-leaf deletion removes every row and column.
\end{enumerate}
Whenever these conditions hold, the core is a three-ball.  Condition~3 is
decidable in $O(k+f+I)$ structural operations after source preparation,
where $I$ is the number of nonzero source incidences.  A successful trace
has a verifier with the same linear bound and
$O((k+f)\log(k+f+1))$-bit labels.
\end{theorem}

\begin{proof}
A successful sequence in condition~1 ends in the empty matrix.  Apply
Lemma~\ref{geo:lem-permanent-pivot} backwards through the sequence to obtain
condition~2.

Suppose condition~2 holds.  Let its unique edge-labelled perfect matching
assign one column to each row.  Every matched entry equals $1$, since two
parallel edges in such an entry would give two matchings.  Name each row
by its matched column, and draw a directed edge $x\to y$ whenever $y\ne x$
and that row has a nonzero entry in column $y$.  A directed cycle would
allow the matched columns along the cycle to be exchanged, producing a
second perfect matching.  The directed graph is therefore acyclic.  A
sink row has no nonzero entry except its unit matched entry, and is a
source row leaf.  Deleting it and its column leaves a matrix whose
permanent is still $1$.  Induction proves condition~3.

Finally, a row leaf has zero in every other live column.  Thus the fill
term in \eqref{geo:eq-positive-pivot} vanishes for all surviving entries.
Every source row-leaf deletion is consequently a valid positive singleton
pivot, and condition~3 implies condition~1.  Lemma~\ref{geo:lem-slide}
realizes these pivots geometrically.  Corollary~\ref{geo:cor-core-handlebody}
gives genus $k-f=0$.

For discovery, store each row's live support size and an incidence list
from each column to its source rows.  Put the eligible rows in a queue.
On deleting a column, update its incident rows; each source incidence is
processed a bounded number of times.  A row enters the queue only when its
live support becomes one, with stale entries discarded if its support
subsequently becomes zero.  This gives the stated structural bound.

An independent checker need not recreate a sequence of mutable matrices.
Assign each deleted column its deletion time.  For every supplied pivot,
check that its source entry equals $1$, and that every other source column
in that row has an earlier deletion time.  Check that each row and column
appears exactly once.  These source conditions state precisely that each
step is a row-leaf deletion.  Scanning the immutable incidences verifies
all conditions in linear structural work.  Identifiers, times and counters
have logarithmic bit length.
\end{proof}

\begin{remark}[Scope of completeness]
\label{geo:rem-class-completeness}
The theorem is complete for the specified raw geometric cancellation
model on the supplied handle structure.  It is not a recognition theorem
for all ball cores.  A ball can require a different disk system, more
general slides, a different spanning tree, or other geometric operations
before this model succeeds.  Even when the source permanent is one, a
poor arbitrary singleton pivot can create positive fill and become stuck.
The source row-leaf procedure avoids this order dependence.  The permanent
is used only in the proof and is never computed by the algorithm.
\end{remark}

\begin{proposition}[A bounded-genus extension]
\label{geo:prop-fixed-genus}
Let a connected authenticated core have $k$ non-tree one-handles, $f$
two-handles, no three-handles, and $g=k-f\ge0$.  Fix a set $S$ of $g$
one-handles to survive.  A complete raw singleton cancellation with these
survivors exists if and only if the square source matrix on all rows and
on columns outside $S$ has permanent one, equivalently if source row-leaf
deletion succeeds when the columns in $S$ are treated as already known.
Therefore exhaustive search over survivor sets is complete within this
raw model in
\[
 \binom{k}{g}\,\operatorname{poly}(T)
 \tag{G5}\label{geo:eq-genus-search}
\]
work when $k+f=O(T)$.  This is polynomial for fixed $g$ and
$2^{O((\log(T+1))^2)}\operatorname{poly}(T)$ when $g=O(\log(T+1))$.
\end{proposition}

\begin{proof}
Restrict any successful positive-pivot sequence to the non-survivor
columns.  The restricted square matrix follows the same positive-pivot
rule and ends empty, so its permanent is one by
Lemma~\ref{geo:lem-permanent-pivot}.  Conversely, apply the row-leaf order
of Theorem~\ref{geo:thm-ball-closure} to this restricted matrix.  A donor
row can have arbitrary source entries in $S$.  Previously eliminated
columns have images supported only in $S$ and still earlier columns.
Consequently they introduce no occurrence of a different live
non-survivor column.  The unique live non-survivor occurrence remains
exactly one, and every cancellation is geometrically valid by
Lemma~\ref{geo:lem-slide}.  All $f=k-g$ rows are consumed.  Enumerate the
$\binom{k}{g}$ possible survivor sets and apply the source test to each.
Finally $\binom{k}{g}\le k^g$ gives the displayed asymptotic consequences.
\end{proof}

The exhaustive survivor search in this proposition is a theoretical
extension of the core certificate, not an assertion that the delivered
implementation searches every such set.  Its geometric use is new to this
consumer; the corresponding immutable-source closure principle was
already proved in report~56 \cite{ProveIt2026}.  It also does not bound a
hierarchy's genera or prove that every unknot enters this cancellation
class.

\subsection{Interval bundles and the actual marked interface}
\label{geo:bundles}

\begin{lemma}[An orientable interval bundle over a bordered surface]
\label{geo:lem-ibundle}
Let $F$ be a connected compact surface with nonempty boundary, and let
$p:P\to F$ be an interval bundle whose total space is orientable.  Then
$P$ is a handlebody of genus
\[
 1-\chi(F).
 \tag{G6}\label{geo:eq-bundle-genus}
\]
If $F$ is orientable of genus $h$ with $b>0$ boundary components, this genus
is $2h+b-1$.  If $F$ is nonorientable with $q$ crosscaps and $b>0$ boundary
components, it is $q+b-1$.
\end{lemma}

\begin{proof}
A connected compact surface with nonempty boundary admits
$1-\chi(F)$ disjoint properly embedded arcs that cut it into a disk.  For
example, take the co-core arcs in a decomposition of $F$ into one disk and
$1-\chi(F)$ surface one-handles.  The inverse image of each arc under $p$
is a properly embedded disk.  Cutting $P$ along these disks gives the
interval bundle over a disk, namely $D^2\times I$, a three-ball.  The disks
are therefore a complete meridian system.  Orientability of the total
space ensures that reconstructing it uses orientable three-dimensional
one-handles, proving the handlebody assertion and the genus formula.  The
two displayed specializations follow from surface Euler characteristic.
\end{proof}

In particular, the orientable interval bundle over a M\"obius band is a
solid torus.  Requiring the base to be orientable would incorrectly discard
this useful case.  Conversely, the nonempty-boundary hypothesis must be
checked: for example, $S^2\times I$ is not a handlebody, and the arc-cutting
argument does not apply to a closed base.

\begin{lemma}[Maximal marked runs give the disk interfaces]
\label{geo:lem-interface}
Assume the full existing source-parallelity conclusion: in the cut
three-manifold the selected total interval bundle $p:P\to F$ is embedded,
its complementary inherited core is $K$, and
\[
 P\cap K=p^{-1}(A),
 \qquad A\subseteq\partial F,
 \tag{G7}\label{geo:eq-intersection}
\]
where $A$ is the union of the closed core-labelled base arcs, with all
source endpoint contacts included.  On each completed base boundary
circle, amalgamate consecutive core-labelled arcs into maximal connected
runs.  If no run is a whole boundary circle, then every component of
$P\cap K$ is an embedded disk.  Distinct such components have separated
closures and admit compatible corner collars.  If a run is a whole
boundary circle, its inverse image is an annulus and is outside this
disk-interface conclusion.
\end{lemma}

\begin{proof}
The completed base boundary is a compact one-manifold.  A finite union of
its closed cell arcs has connected components that are either closed
intervals or entire circles; an interval is maximal precisely after all
consecutive selected arcs have been amalgamated.  Distinct interval
components are separated by nonempty open complementary arcs.  An
interval bundle over a closed interval is a disk.  Over a boundary circle
of a surface, an interval bundle with orientable total space restricts to
an annulus: the base has an annular collar there, so orientation of the
total space rules out reversal of the fiber around that circle.

Equation~\eqref{geo:eq-intersection} is a statement about the embedded
bundle and core, not merely their quotient data.  Thus each inverse-image
disk is embedded also on the core side.  Distinct components are disjoint
compact subsets of this embedded product boundary; there are only finitely
many core arcs, and hence finitely many such components.  Choose disjoint
small neighborhoods and use the fiber charts and boundary collars to
round their corners compatibly.  The rounded disks still have disjoint
closures.  An apparent shared corner between different bundle components,
or an alternating corner that was not identified in the completed base,
would contradict the assumed embedded interval-bundle structure and its
full product-intersection formula.  This is why the source theorem's edge
and vertex contacts are needed.
\end{proof}

The existing geometric theorem supplying
\eqref{geo:eq-intersection} is the trusted mathematical input here; its
source charts are reconstructed by the geometry verifier
\cite{ProveIt2026}.  The new ordered boundary word certifies which maximal
components of $A$ are intervals and which are circles.  A cyclic word scan
alone would not establish \eqref{geo:eq-intersection}.  In particular, a
single face port is not automatically one disk interface.  Several ports
can constitute one disk, or together form an annulus.  Independently
shrinking every face port would change the attachment and is not a valid
substitute for using its true maximal run.

Each maximal run has one owner component of the core, as verified from the
source attachment charts.  This owner is a core \emph{component}, not
necessarily one core handle: consecutive arcs can face distinct handles
within that same component.

\subsection{Disk gluing gives an exact genus formula}
\label{geo:disk-gluing}

\begin{theorem}[Handlebodies glued along separated disks]
\label{geo:thm-disk-gluing}
Let $\{H_v:v\in V\}$ be finitely many connected orientable handlebodies of
genera $g_v$.  Specify pairs of embedded closed disks on their boundaries,
all with pairwise disjoint closures before gluing, and a homeomorphism for
each pair.  Suppose that the quotient, with corners rounded, is orientable.
Let $\Gamma$ be the incidence multigraph, with one vertex per handlebody
and one edge per disk pair; loops and parallel edges are allowed.  Every
connected component $C$ of $\Gamma$ gives a handlebody component of the
quotient, of genus
\[
 G_C=\sum_{v\in V(C)}g_v+|E(C)|-|V(C)|+1.
 \tag{G8}\label{geo:eq-gluing-genus}
\]
It is a three-ball if and only if every constituent genus is zero and
$C$ is a tree.
\end{theorem}

\begin{proof}
The pieces meet only in the specified disk identifications, so the graph
records connectedness exactly.  Work with one connected graph component
and choose a spanning tree.

Choose a complete meridian system in each $H_v$ that misses all its marked
disks.  Such a choice always exists: first move the finitely many marked
boundary disks, by a boundary isotopy, into small disjoint neighborhoods
away from any chosen meridian boundaries; extend the isotopy over a collar
of the handlebody and pull the meridian system back.  Thus this step does
not impose an extra condition on the placement or gluing maps of the
marked disks.

Cut the quotient along all these local meridian disks.  Also cut along the
gluing disk corresponding to each edge outside the chosen spanning tree.
The remaining manifold is a tree of balls glued along boundary disks.  It
is a ball, by repeatedly deleting a leaf ball and applying the elementary
ball-attachment lemma; compare \cite[Lemma~1.3]{Hatcher2007}.  Consequently
the local meridian disks together with the non-tree gluing disks form a
complete meridian system.  Their total number is
\[
 \sum_v g_v+|E|-|V|+1.
\]
The orientability hypothesis yields an orientable handlebody of this
genus.  Every $g_v$ is nonnegative, and $|E|-|V|+1$ is the nonnegative cycle
rank of a connected graph.  Their sum is zero exactly in the stated ball
case.
\end{proof}

\begin{corollary}[A terminal classifier for a marked normal cut]
\label{geo:cor-normal-cut-classifier}
Take the existing authenticated normal-cut parallelity decomposition.
Assume that every core component has a successful certificate from
Corollary~\ref{geo:cor-core-handlebody}, every interval-bundle base has
nonempty boundary, and every maximal core-labelled boundary run is proper.
Then each component of the entire cut is a handlebody.  Its genus is given
by \eqref{geo:eq-gluing-genus}, using the core genera from
\eqref{geo:eq-core-genus}, the bundle genera from
\eqref{geo:eq-bundle-genus}, and the actual maximal-run incidence graph.
In particular, the entire cut is one three-ball exactly when that graph is
connected and acyclic and all its vertex genera are zero.
\end{corollary}

\begin{proof}
The core certificates give the handlebody pieces.  Lemma~\ref{geo:lem-ibundle}
gives the same conclusion for the bundle pieces, including those with
nonorientable bases.  Lemma~\ref{geo:lem-interface} supplies their separated
disk interfaces.  Apply Theorem~\ref{geo:thm-disk-gluing}.
\end{proof}

This classification uses only the existence of the homeomorphisms supplied
by the core cancellations.  Such a homeomorphism automatically transports
the originally authenticated boundary disks to disjoint embedded disks.
Because Theorem~\ref{geo:thm-disk-gluing} allows arbitrary disk placements,
one need not explicitly transport their parametrizations to compute the
genus.  Explicit marking transport would again be necessary for a later
operation that needs specified curves or disks, rather than this terminal
topological type.

An implementation must also preserve the distinction between an actual
piece and a multiplicity-weighted profile of pieces.  A bundle component
with a core port has a source owner and multiplicity one.  Unattached
repeated profiles can contribute many isolated graph components.  The
one-ball conclusion requires the total connected-component count to be
one, including those unattached components.

\subsection{A concrete extension: absorbing disc-base annuli}
\label{geo:annulus-absorption}

A full core-labelled boundary circle prevents disk gluing.  For a disk base,
however, its bundle is a two-handle.  Lackenby's refined algorithm already uses
this replacement for interior parallelity bundles\cite[Section~10, step~$(2')$]{Lackenby2026}.
The contribution below supplies the source-chart incidence row and proves
compatibility with the present geometric cancellation certificates.

\begin{proposition}[Source-aware absorption of a disc-base bundle]
\label{geo:prop-annulus-absorption}
In the authenticated source-parallelity decomposition, let $P\to F$ be a
connected bundle component whose base is a disk and whose entire base
boundary is core-labelled.  Then $P$ can be absorbed into its owner core
component as one two-handle.  No zero-, one-, or three-handles need be
added.

Use the full ordered boundary port circle and its native source charts.
Every port of core dimension one is a face-sector strip running from one
end of that source one-handle to the other, and every port of core dimension
zero lies in a zero-handle boundary.  Thus the new attaching circle has
raw unsigned incidence
\[
 n^{\mathrm{new}}_h
 =\#\{\text{ports on this circle of core dimension one and owner }h\}.
 \tag{G9}\label{geo:eq-absorbed-row}
\]
After deleting spanning-forest columns, this row may be appended to the
existing authenticated core incidence matrix.  The geometric pivot,
saturation, and source row-leaf theorems then apply to the augmented core.
The assertion holds simultaneously for any collection of such bundle
components.  There are $O(T)$ possible absorbed components and $O(T)$
source port occurrences in all their new rows together.
\end{proposition}

\begin{proof}
An interval bundle over a disk is $D^2\times I$.  Its intersection with the
core is, by the full source product-intersection conclusion,
\[
 p^{-1}(\partial F)=\partial D^2\times I.
\]
This is exactly the attaching region of a three-dimensional two-handle.
The annulus is connected, so its image lies in one core component.  The
two horizontal disks are the free boundary of the attached handle.

To authenticate its attaching circle relative to the core's zero/one
structure, choose the midpoint section of its fibers around $\partial F$.
The native fiber transition maps preserve this midpoint, including maps
that reverse the fiber.  The resulting curve is therefore an embedded
circle in the attaching annulus.  On a core-zero port it lies in the
boundary of the indicated zero-handle.  The native constructor's
core-one ports come precisely from the \texttt{face\_sector} sides of
selected edge-cap polygons.  Each such side traverses the lateral boundary
of its core one-handle between the two source tetrahedron ends, even when
those ends belong to the same tetrahedron or zero-handle.  The midpoint
curve consequently meets that one-handle's co-core exactly once.  The
source face flags, rather than equality of the tetrahedron names, identify
these ends.  This proves \eqref{geo:eq-absorbed-row}.

There is also an explicit signed chart check, although the saturated
classifier needs only unsigned occurrences.  Let the oriented source ends
of a core one-handle $h$ have flags $(t_h,f_h)$ and $(u_h,g_h)$.  For the
face-sector port $p$, inspect the initial and terminal cap-edge endpoint
flags in its local direction and set
\[
 \delta_p=
 \begin{cases}
 +1,&\text{initial}=(t_h,f_h),\quad\text{terminal}=(u_h,g_h),\\
 -1,&\text{initial}=(u_h,g_h),\quad\text{terminal}=(t_h,f_h).
 \end{cases}
 \tag{G10}\label{geo:eq-port-sign}
\]
Require one of these two cases.  If the port is traversed with sign $s_p$
in the canonical boundary word, its raw letter is $x_h^{s_p\delta_p}$.
A core-zero port stays in its owner zero-handle, while a core-one port
moves between the two stated ends.  The resulting sequence must be a
closed path in the source zero/one graph.  Both tetrahedron and face flags
are needed: equality of tetrahedron labels alone does not distinguish the
two ends of a self-pairing.  No abelian signed sum replaces the unsigned
count in \eqref{geo:eq-absorbed-row}.

These port charts lie in the boundary of the core zero/one-handle union
outside the attaching regions of the original core two-handles.  This is
part of the genuine source dual-handle structure: the complete lower
attachments of complementary handles stay in the core, while the selected
product contacts occupy their recorded, disjoint boundary regions.
Compatible small collars make the new midpoint circles disjoint from the
original two-handle attaching curves.  Hence all the new handles can be
placed at the same two-handle stage, with exactly the authenticated rows
just computed.  Tree contraction deletes only the tree entries, as before.
The preceding geometric cancellation theorems apply to this augmented
handle structure without introducing a three-handle.

Different bundle components are disjoint in the embedded source bundle,
so their attaching annuli and midpoint curves can be regularized
disjointly.  Absorbing them leaves the marked interfaces of the remaining
components intact.  Each absorbed component contains at least one core
port, and every port has one actual owner.  The native $O(T)$ port bound
therefore bounds both the number of new two-handles and the total number
of port occurrences used in their rows.
\end{proof}

This proposition is a proved extension of the conservative classifier;
absorption is not performed by the delivered disk-interface consumer.
It suggests a direct next implementation: identify the certified base
signature $(\chi,b)=(1,1)$ with its complete boundary core-labelled,
append its source-authenticated incidence row, run the core cancellation
checker again, and apply disk gluing to the remaining bundle pieces.
The base's surface structure and the full port charts must be verified;
the two numerical values $(\chi,b)$ alone would not supply the necessary
geometric conclusion.  A bundle over an annulus or a M\"obius band is a
solid torus rather than a two-handle, so the proposition does not silently
extend to those cases.

\subsection{Why one ball after a connected cut certifies the unknot}
\label{geo:terminal}

\begin{proposition}[The one-ball cut criterion]
\label{geo:prop-terminal}
Let $M$ be a connected compact orientable three-manifold with boundary one
torus.  Let $S\subset M$ be a connected orientable properly embedded
surface with nonempty boundary.  If the manifold obtained by cutting $M$
along $S$ is one three-ball, then $S$ is a disk and $M$ is a solid torus.
If $M$ is an authenticated knot exterior in $S^3$, this certifies that the
knot is the unknot.
\end{proposition}

\begin{proof}
The ambient and surface orientability imply that $S$ is two-sided.  Let
$M_S$ denote the manifold cut along a product neighborhood of $S$.
Euler characteristic under this cut gives
\[
 \chi(M_S)=\chi(M)+\chi(S).
\]
For a compact orientable three-manifold,
$\chi(M)=\tfrac12\chi(\partial M)$; this also follows by doubling and using
the vanishing Euler characteristic of a closed odd-dimensional manifold.
The torus boundary gives $\chi(M)=0$, while the ball hypothesis gives
$\chi(M_S)=1$.  Hence $\chi(S)=1$.  Classification of connected orientable
surfaces with nonempty boundary implies that $S$ is a disk.

Reconstructing $M$ from the cut ball inserts a neighborhood of this disk,
namely one three-dimensional one-handle attached along its two boundary
copies.  Orientability makes the result a solid torus.

For a knot exterior, an essential meridian disk of this exterior has
boundary equal, up to orientation, to the knot's preferred longitude:
that primitive slope generates the kernel of the map from the boundary
torus's first homology to the exterior's first homology.  Inside the knot's
tubular neighborhood this longitude cobounds an annulus with the knot.
The disk and annulus form an embedded spanning disk for the knot, proving
that it is the unknot.
\end{proof}

The hypotheses must be used literally.  Finding a ball among several cut
components is insufficient.  A cut component of positive handlebody genus
is insufficient.  A multiple normal vector can describe a disconnected
surface, so its connectedness and orientability cannot be inferred from a
successful core certificate.

\subsection{An all-size Fibonacci family with ball cores}
\label{geo:fibonacci}

We now prove a uniform success theorem for the native complementary core
of the layered solid-torus fixture from report~67
\cite{ProveIt2026}.  This gives a construction and proof for every family
member; finite timings play no role in the theorem.

Let $F_1=F_2=1$ and $F_{j+2}=F_{j+1}+F_j$.  For $t\ge1$, number the
tetrahedra $0,\ldots,t-1$.  For each $j<t-1$, glue face $0$ of tetrahedron
$j$ to the next tetrahedron using the vertex permutation
\[
 P_0=(2,1,3,0),
\]
and face $1$ using
\[
 P_1=(0,3,1,2).
\]
At the last tetrahedron, pair face $0$ to face $1$ using
$P=(1,2,3,0)$.  Thus the unpaired faces are the top faces $2$ and $3$.
For an arbitrary positive integer $m$, use normal-coordinate rows
\[
 m\bigl(F_{t-j+1},F_{t-j+1},0,0,0,0,F_{t-j}\bigr),
 \qquad 0\le j<t.
 \tag{G11}\label{geo:eq-fibonacci-vector}
\]
The primitive fixture is the standard meridian of this layered solid
torus.  The following theorem concerns its native \emph{complementary
core}, for every positive scale, not the number of components of the
whole cut.

\begin{theorem}[Uniform ball-core certificates for the Fibonacci fixtures]
\label{geo:thm-fibonacci}
For every $t\ge1$ and every positive integer $m$, the native core of the
cut specified by \eqref{geo:eq-fibonacci-vector} has handle counts
\[
 h_0=4t,\qquad h_1=6t-3,\qquad h_2=2t-2,\qquad h_3=0,
 \tag{G12}\label{geo:eq-fibonacci-counts}
\]
and is a three-ball.  With the native source ordering described below,
the increasing-identifier spanning tree followed by immutable-source
row-leaf deletion removes every two-handle.  The certificate has
$O(t\log(t+1))$ bits and $O(t)$ source incidences, independently of $m$ and
of the numerical normal weights.
\end{theorem}

\begin{proof}
We describe the zero/one-handle graph, enumerate all interior edge links,
and exhibit a row-leaf order valid for every $t$.

\paragraph{The zero/one-handle graph.}
At tetrahedron $j$ the occupied triangle types are $0$ and $1$, and the
occupied quadrilateral type is column $6$.  The four exceptional chambers
are denoted $A_j,B_j,U_j,V_j$: the two quad-end chambers, followed by the
triangle tips at vertices $0$ and $1$.  In this order their identifiers are
$4j,4j+1,4j+2,4j+3$.  Applying the two layer permutations to the central
face region and its occupied corner regions gives the following
one-handles in each nonfinal layer:
\[
\begin{array}{c|c}
\text{identifier}&\text{ends}\\ \hline
6j   &(B_j,A_{j+1})\\
6j+1 &(V_j,V_{j+1})\\
6j+2 &(A_j,U_{j+1})\\
6j+3 &(A_j,B_{j+1})\\
6j+4 &(U_j,U_{j+1})\\
6j+5 &(B_j,V_{j+1}).
\end{array}
 \tag{G13}\label{geo:eq-fibonacci-graph}
\]
The last self-pairing gives three more one-handles, of identifiers
$6t-6,6t-5,6t-4$, with ends
\[
 (B_{t-1},A_{t-1}),\qquad
 (V_{t-1},B_{t-1}),\qquad
 (A_{t-1},U_{t-1}),
\]
respectively.  These lists already prove $h_0=4t$ and $h_1=6t-3$.

For $t\ge3$, the increasing-identifier spanning tree consists of
\[
 \{0,\ldots,10\}\ \cup\!
 \bigcup_{j=2}^{t-2}\{6j,6j+1,6j+2,6j+3\}.
 \tag{G14}\label{geo:eq-fibonacci-tree}
\]
Indeed, the first six edges in \eqref{geo:eq-fibonacci-graph} form two
trees on four vertices each.  Edges $6,7,8,9$ attach the four vertices of
tetrahedron $2$, and edge $10$ joins the two previous trees.  Thus the first
eleven edges form a tree on the first twelve vertices.  Thereafter the
first four edges of a layer each attach one new vertex to the existing tree.  The other two
edges of that layer join vertices already connected, as do the three
final self-pairing edges.  Thus \eqref{geo:eq-fibonacci-tree} has exactly
$4t-1$ edges and is the claimed spanning tree.  When $t=2$ the tree is
$\{0,\ldots,6\}$, and when $t=1$ all three one-handles form the tree.

The bottom self-pairing identifies vertices $1,2,3,0$ successively.  The
layer permutations propagate this identification to every preceding
tetrahedron.  All source vertices consequently lie in the single boundary
vertex class containing the top vertices.  There are no interior vertex
links and hence no core three-handles.

\paragraph{The complete edge inventory.}
Write $(j;ab)$ for local edge $\{a,b\}$ in tetrahedron $j$.  For every
$0\le j\le t-2$ there is an interior edge $E_j$ represented by $(j;23)$.
For $0\le j\le t-4$, its cyclic list of local occurrences is
\[
 (j;23),\ (j+1;12),\ (j+2;13),\ (j+3;01),\
 (j+2;02),\ (j+1;03).
 \tag{G15}\label{geo:eq-fibonacci-link-general}
\]
The exit faces are, in order, $1,0,0,3,3,2$.  Each transition follows by
applying the indicated layer permutation or its inverse.  For the
penultimate edge, $j=t-3$, the bottom self-pairing shortens the list to
\[
 (j;23),\ (j+1;12),\ (j+2;13),\ (j+2;02),\ (j+1;03),
 \tag{G16}\label{geo:eq-fibonacci-link-penultimate}
\]
and for the terminal edge, $j=t-2$, the list is
\[
 (j;23),\ (j+1;12),\ (j+1;23),\ (j+1;03).
 \tag{G17}\label{geo:eq-fibonacci-link-last}
\]
These cycles are disjoint as sets of local edge occurrences.  For
$t\ge3$, the remaining occurrences form exactly the following three
boundary classes:
\[
\begin{split}
 &\{(0;01)\},\\
 &\{(0;02),(0;13),(1;01)\},\\
 &\{(0;03),(0;12),(1;02),(1;13),(2;01)\}.
\end{split}
 \tag{G18}\label{geo:eq-fibonacci-boundary-edges}
\]
Their transitions again follow directly from the face maps; each class
meets a top boundary face.  The interior lists use $6t-9$ local edge
occurrences and the boundary lists use the remaining nine, so this is the
complete edge inventory.  For $t=2$, omit $(2;01)$ from the last boundary
class; the sole interior cycle is \eqref{geo:eq-fibonacci-link-last}.  For
$t=1$, the three boundary classes in the only tetrahedron are
$\{01\}$, $\{02,13\}$ and $\{03,12,23\}$, and there is no interior edge.

Every $E_j$ has positive normal edge weight, equal to $mF_{t-j}$.  The
native core therefore has exactly two two-handles around it, at rank zero
and at the full edge weight.  Thus there are exactly $2t-2$ core
two-handles, proving the remaining count in
\eqref{geo:eq-fibonacci-counts}.

\paragraph{The literal attaching rows.}
Give these two handles identifiers $2j$ and $2j+1$.  Reading their endpoint
ranks around \eqref{geo:eq-fibonacci-link-general} and using
\eqref{geo:eq-fibonacci-graph} gives the following one-handle crossing
lists for $0\le j\le t-4$:
\[
\begin{split}
 2j:\quad &(6j+5,\ 6j+7,\ 6j+13,\ 6j+17,\ 6j+9,\ 6j),\\
 2j+1:\quad &(6j+3,\ 6j+6,\ 6j+14,\ 6j+16,\ 6j+10,\ 6j+2).
\end{split}
 \tag{G19}\label{geo:eq-fibonacci-full-general}
\]
For $j=t-3$, the shortened link gives
\[
\begin{split}
 2j:\quad &(6j+5,\ 6j+7,\ 6j+13,\ 6j+9,\ 6j),\\
 2j+1:\quad &(6j+3,\ 6j+6,\ 6j+14,\ 6j+10,\ 6j+2),
\end{split}
 \tag{G20}\label{geo:eq-fibonacci-full-penultimate}
\]
and for $j=t-2$ it gives
\[
\begin{split}
 2j:\quad &(6j+5,\ 6j+7,\ 6j+6,\ 6j),\\
 2j+1:\quad &(6j+3,\ 6j+6,\ 6j+8,\ 6j+2).
\end{split}
 \tag{G21}\label{geo:eq-fibonacci-full-last}
\]
For example, the two endpoint ranks at a face arm select its corner-end
and central exceptional regions respectively; applying the face
permutation selects the corresponding endpoint in
\eqref{geo:eq-fibonacci-graph}.  The cyclic edge lists enumerate every
such passage, so there are no additional crossings.  All entries in the
displayed lists have multiplicity one.  The numerical values of the full
edge ranks do not change which exceptional endpoint is selected.

\paragraph{A uniform source row-leaf order.}
Assume $t\ge3$.  Name the non-tree columns
\[
 \begin{array}{ll}
 a_j=6j+5,&1\le j\le t-2,\\
 b_j=6j+4,&2\le j\le t-2,\\
 x=6t-6,\quad y=6t-5,\quad z=6t-4.&
 \end{array}
 \tag{G22}\label{geo:eq-fibonacci-columns}
\]
Use $a_0$ and $b_1$ as formal names for absent columns: their identifiers
$5$ and $10$ belong to the spanning tree.  Deleting the tree entries from
\eqref{geo:eq-fibonacci-full-general}--\eqref{geo:eq-fibonacci-full-last}
gives the following support table, where absent formal columns are
omitted:
\[
\begin{array}{c|c|c}
\text{range}&\text{row }2j&\text{row }2j+1\\ \hline
0\le j\le t-4 &\{a_j,a_{j+2}\}&\{b_{j+1},b_{j+2}\}\\
j=t-3 &\{a_{t-3},y\}&\{b_{t-2},z\}\\
j=t-2 &\{a_{t-2},x,y\}&\{x,z\}.
\end{array}
 \tag{G23}\label{geo:eq-fibonacci-reduced}
\]
This is a square matrix of size $2t-2$.

First peel the odd rows in increasing order.  Since $b_1$ is absent, the
successive deleted columns are
\[
 b_2,b_3,\ldots,b_{t-2},z,x.
\]
The initial range is empty when $t=3$; the same rule still deletes $z$
and then $x$.

After $x$ is removed, view the even rows as edges on the formal vertices
$a_0,a_1,\ldots,a_{t-2},y$.  The first range of rows supplies the parity-step
edges $a_j a_{j+2}$.  The last two even rows join the final vertices of
these two parity chains to $y$.  Their union is therefore a single path,
with one endpoint the absent column $a_0$.  Peel along this path from that
endpoint.  The deleted columns are the even-indexed $a$'s in increasing
order, followed by $y$, followed by the odd-indexed $a$'s in decreasing
order, ending at $a_1$.  Every row is a source row leaf when it is used.
Thus all even rows are removed as well.

For $t=2$, the two surviving columns have identifiers $7$ and $8$, and the
two reduced rows form the identity matrix.  For $t=1$, the reduced matrix
is empty.  Theorem~\ref{geo:thm-ball-closure} now proves that the core is a
ball for every $t$.

Finally, multiplying \eqref{geo:eq-fibonacci-vector} by a positive integer
changes neither its occupied disk types nor the positivity of any relevant
edge weight.  Every exceptional chamber, endpoint rank, source crossing
list and leaf order above therefore remains unchanged for all $m\ge1$.
There are $O(t)$ forest edges and pivots, and each crossing list has bounded
length, proving the certificate-size assertion.
\end{proof}

The primitive coordinates in this family already have exponential
numerical weight in $t$, whereas the core certificate is linear in its
number of source records.  The theorem establishes an exact, nontrivial
independence from normal multiplicity for this core decision.  It does not
identify a coordinate's numerical magnitude with its binary input length,
and it does not extend the one-ball cut criterion to disconnected
multiples.  The marked interval-bundle assembly remains necessary to
classify the entire cut.
